Eight PID tunings were compared on a three-state exothermic CSTR whose cooling valve is limited in amplitude and in slew rate, under a single protocol with true-temperature metrics. Four conclusions follow.
\begin{enumerate}
    \item \textbf{The anti-windup implementation can decide the verdict on a tuning rule.} With directional clamping that accounts for amplitude and rate saturation and for the sign of the reverse-acting loop, every tuning tracks the nominal step; with the textbook amplitude-only rule, ZN and CC lose temperature control ($\mathrm{IAE}$ of $104.123$ and $102.668$ instead of $0.940$ and $1.017$~K$\cdot$min) and the high-gain modern tunings overshoot several times more. Comparisons that report the failure of empirical rules should state their anti-windup logic.
    \item \textbf{Robustness, not stability, separates the tunings.} The small margins of ZN and CC ($M_s = 4.14$ and $4.54$ on the design model) are masked on a clean step by the relative-degree-two plant and the rate limit, but under sensor noise they produce the largest thermal deviation, about twice that of SIMC at $\sigma = 1$~K. The FOPTD design model is a conservative proxy for the peak sensitivity of the real loop but can overestimate the phase margin by up to $5.4^\circ$.
    \item \textbf{Constrained optimization and setpoint weighting give the best servo compromise.} Multi-start ITAE optimization with $M_s \le 1.6$ on the dead-time loop yields the fastest robust PID ($\mathrm{IAE} = 0.827$~K$\cdot$min, constraint active), while a single-start search stops at a PI-like local optimum. Setpoint weighting reduces overshoot and valve travel for AMIGO and Optimal ITAE without changing their disturbance response, and the SIMC locus shows that tracking speed is bought with robustness along a front without interior optimum.
    \item \textbf{Tracking-error variance is governed by physics.} Under simultaneous uncertainty the IAE variance of every tuning is explained mainly by feed concentration and heat-transfer capacity, i.e. by the heat-removal limit of the reactor, whereas the valve activity of low-margin loops is driven by kinetics and sensor dead time.
\end{enumerate}
The released code reproduces every table and figure and accepts new controllers behind the same interface, so that future tuning methods, and nonlinear or learning-based controllers, can be compared with these robust PID baselines under identical plant, actuator and protocol. Future work will extend the analysis to gain-scheduled PID and to operating points inside the open-loop unstable region.
